% Source: https://artofproblemsolving.com/community/c3h1573199_need_help_on_mathcounts_geometry_problems
% Stable identity: aops-topic-1573199-q01
\begin{asy}
usepackage("amsmath");
size(5cm);
defaultpen(fontsize(10pt));
real s=28,base=20; pair B=(0,0),C=(base,0),A=(base/2,sqrt(s^2-(base/2)^2));
real t=0.35; pair Ee=B+t*(C-B);
pair D=B+t*(A-B), F=C+(1-t)*(A-C);
fill(A--D--Ee--F--cycle,paleyellow);
draw(A--B--C--cycle,linewidth(0.9));draw(D--Ee--F,blue+linewidth(0.9));
dot(D);dot(Ee);dot(F);
label("$A$",A,N);label("$B$",B,SW);label("$C$",C,SE);label("$D$",D,W);label("$E$",Ee,S);label("$F$",F,E);
label("$28$",(A+B)/2+(-1.2,1.8),W);label("$28$",(A+C)/2+(1.2,1.8),E);draw((0,-3.8)--(20,-3.8),gray,Arrows(TeXHead));label("$20$",(10,-3.8),S);
label("\mbox{}",(min(currentpicture,user=true)+max(currentpicture,user=true))/2,black); // keeps the text after the figure black when LaTeX typesets the labels inline
\end{asy}

\textbf{$DB=DE$.} Since $DE\parallel AC$, corresponding angles give $\angle DEB=\angle ACB$. As $AB=AC$, $\angle ACB=\angle ABC=\angle DBE$. So triangle $DBE$ has equal angles at $B$ and $E$, hence $DB=DE$.

\textbf{Perimeter.} In parallelogram $ADEF$, opposite sides are equal: $AF=DE$ and $EF=AD$. Therefore
\[
\text{perimeter}=2(AD+DE)=2(AD+DB)=2AB=2\cdot28=\boxed{56}.
\]
The answer does not depend on where $E$ is on $BC$; the length $BC=20$ is not needed.
